% !TEX root =./ECE448_Practice_main.tex
%
% L5 practice set (GNU Radio Introduction). Short by design: the L5 Lab packet
% carries the hands-on work. Blank copy: \begin{solution}[h] reserves h of work
% space. SOLUTIONS copy (\iskey=1) prints the red worked solutions.

\fullwidth{\textbf{Learning Objective 2.3: } I can install GNU Radio and use GNU Radio
Companion (GRC) to build, inspect, and run a basic flowgraph.
\\[2pt]\textbf{Learning Objective 2.4: } I can capture a real RF signal using an SDR
and interpret its time- and frequency-domain representations.
\\[4pt]\normalfont\small Show your work. Write a \textbf{Documentation} line at the
end (\textbf{None} if you did not collaborate).}

\begin{questions}

%%%%%%%%%%%%%%%%%%%%%%%%%%%%%%%%%%%%%%%%%
\question \textbf{Reading a flowgraph.} A classmate's flowgraph has a Signal Source,
a Throttle, and a QT GUI Frequency Sink. GRC draws the arrow from the Throttle to
the sink in red and will not run.
\begin{parts}

	\part The source and throttle ports are blue and the sink's input is orange.
	What is wrong, and what are two ways to fix it?
		\begin{solution}[1.0in]
		The source and throttle carry complex samples (blue) and the sink expects
		floats (orange); connected ports must have the same type. Either set the
		sink's type to complex, or insert a conversion block such as Complex to Mag
		or Complex to Real between them.
		\end{solution}

	\part The classmate then replaces the Signal Source with a PlutoSDR Source and
	keeps the Throttle. Explain why the Throttle should come out.
		\begin{solution}[1.0in]
		The PlutoSDR delivers samples at exactly its sample rate, so the hardware already
		paces the flowgraph. A throttle only exists to pace a flowgraph with no hardware,
		and with an SDR it adds a second, competing clock that can drop samples.
		\end{solution}

\end{parts}

\newpage
%%%%%%%%%%%%%%%%%%%%%%%%%%%%%%%%%%%%%%%%%
\question \textbf{Complex sampling and aliasing.} A complex signal source runs at
$f_s = 32$ kHz.
\begin{parts}

	\begin{minipage}[t]{0.58\linewidth}
		\part What range of frequencies can the samples represent?
		\begin{solution}[0.7in]
		\eq{ -\dfrac{f_s}{2} \text{ to } +\dfrac{f_s}{2} = -16 \text{ to } +16\ \text{kHz} }
		\end{solution}
	\end{minipage}\hfill%
	\begin{minipage}[t]{0.37\linewidth}
		\raggedleft \vspace{12pt}
		\ansbox[1.5in]{$-16$ to $+16$ kHz}
	\end{minipage}

	\begin{minipage}[t]{0.58\linewidth}
		\part Where does the frequency sink show a 40 kHz tone, and a $-20$ kHz tone?
		\begin{solution}[1.0in]
		\eq{ 40 - 32 &= 8\ \text{kHz} \\ -20 + 32 &= 12\ \text{kHz} }
		\end{solution}
	\end{minipage}\hfill%
	\begin{minipage}[t]{0.37\linewidth}
		\raggedleft \vspace{12pt}
		40 kHz: \ansbox[1.0in]{$8$ kHz}

		\vspace{4pt}
		$-20$ kHz: \ansbox[1.0in]{$12$ kHz}
	\end{minipage}

	\begin{minipage}[t]{0.58\linewidth}
		\part A PlutoSDR samples at 2.4 MS/s centered on 100.1 MHz. What band does the
		display show? Suppose a strong station at 102.1 MHz leaks past the Pluto's
		filters unattenuated. Where does it appear?
		\begin{solution}[1.2in]
		The display spans $100.1 \pm 1.2$ MHz, 98.9 to 101.3 MHz. The station is
		$+2.0$ MHz from center, outside the $\pm 1.2$ MHz window, so it aliases to
		$2.0 - 2.4 = -0.4$ MHz from center:
		\eq{ 100.1 - 0.4 = 99.7\ \text{MHz} }
		\end{solution}
	\end{minipage}\hfill%
	\begin{minipage}[t]{0.37\linewidth}
		\raggedleft \vspace{12pt}
		band: \ansbox[1.5in]{98.9 to 101.3 MHz}

		\vspace{4pt}
		alias: \ansbox[1.5in]{99.7 MHz}
	\end{minipage}

	\part Explain why complex samples show twice the bandwidth of real samples at the
	same rate.
		\begin{solution}[0.9in]
		Real samples cannot tell a positive frequency from a negative one, so they
		represent only 0 to $f_s/2$, with a mirror image below zero. Complex samples
		record the direction of rotation, so $+f$ and $-f$ are different signals and the
		full $-f_s/2$ to $+f_s/2$ is usable.
		\end{solution}

\end{parts}

\newpage
%%%%%%%%%%%%%%%%%%%%%%%%%%%%%%%%%%%%%%%%%
\question \textbf{Recording and reading IQ.} A File Sink records complex samples.
\begin{parts}

	\begin{minipage}[t]{0.58\linewidth}
		\part Find the size of a 5 s capture at 2.4 MS/s.
		\begin{solution}[0.8in]
		\eq{ 5 \times 2.4\times10^{6} \times 8\ \text{bytes} = 96\ \text{MB} }
		\end{solution}
	\end{minipage}\hfill%
	\begin{minipage}[t]{0.37\linewidth}
		\raggedleft \vspace{12pt}
		\ansbox[1.3in]{96 MB}
	\end{minipage}

	\part You play the capture back through a File Source into a QT GUI Sink, with
	no other blocks. The flowgraph uses a whole processor core and the display
	races. What block is missing, and what do you set it to?
		\begin{solution}[0.9in]
		A Throttle, set to the capture's sample rate, 2.4 MS/s. A File Source has no
		hardware to pace it, so without a throttle it reads samples as fast as the
		laptop can.
		\end{solution}

	\part A teammate sends only the file \texttt{capture.cfile}. Name the two numbers
	you need before the file is useful, and explain why.
		\begin{solution}[0.9in]
		The sample rate and the center frequency. The file holds raw $I$ and $Q$ values
		with no header: without the sample rate the time and frequency axes have no
		scale, and without the center frequency the frequency axis has no reference.
		\end{solution}

\end{parts}

\newpage
%%%%%%%%%%%%%%%%%%%%%%%%%%%%%%%%%%%%%%%%%
\question \textbf{Gain and the DC spike.} A PlutoSDR runs the L5 FM flowgraph,
with the gain set to Manual.
\begin{parts}

	\part With the gain at 0 dB, only the strongest station shows, and the noise
	floor does not move when you raise the gain a little. What is happening, and what
	do you change?
		\begin{solution}[0.9in]
		The gain is too low: weak stations sit below the converter's own noise, which
		sets the floor. Raise the gain until the floor rises a few dB, which shows the
		antenna's noise now sits above the converter's.
		\end{solution}

	\part With the gain at maximum, copies of a strong station appear at frequencies
	where nothing transmits. Name the problem and the fix.
		\begin{solution}[0.9in]
		Overload: the strong station drives the converter past full scale, and the
		clipping creates spurious signals. Lower the gain until the copies disappear,
		keeping it high enough that the noise floor sits a few dB above the
		converter's own.
		\end{solution}

	\part A student tunes the PlutoSDR exactly onto a weak narrowband signal. With DC
	tracking off, the signal hides in a spike at the center of the display; with DC
	tracking on, it disappears. Explain both, and the fix.
		\begin{solution}[1.1in]
		The spike is the receiver's own DC offset, which appears at 0 Hz, the center
		frequency, in every zero-IF SDR. DC tracking subtracts whatever sits at 0 Hz,
		so it removes the offset and the signal with it. Tuning the SDR a little off the
		signal moves the signal away from the center, where neither touches it.
		\end{solution}

\end{parts}

\vspace{1.5em}
\noindent\textbf{Documentation:} \rule[-2pt]{0.6\linewidth}{0.4pt}

\end{questions}
